\documentclass[11pt,a4paper]{article}
\usepackage[utf8]{inputenc}
\usepackage[T1]{fontenc}
\usepackage[spanish]{babel}
\usepackage{amsmath,amsfonts,amssymb}
\usepackage{graphicx}
\usepackage{hyperref}
\usepackage{geometry}
\geometry{margin=2.5cm}
\usepackage{abstract}
\renewcommand{\abstractnamefont}{\normalfont\bfseries}
\usepackage{booktabs}
\usepackage{multirow}
\usepackage{natbib}
\usepackage{parskip}
\usepackage{subcaption}
\usepackage{tikz}

\title{Descomposición Empírica de Modos y Transformación a Grafos del Indice MSCI World: Un Análisis Topológico Multiescala para el Modelado con Redes Neuronales Gráficas}

\author{Nombre y apellidos del autor}

\date{13 de diciembre de 2025 (última revisión a 12:00AM)}

\AtBeginDocument{\renewcommand{\abstractname}{}}

\begin{document}

\maketitle

\begin{abstract}
Este trabajo explora la aplicación de la Descomposición Empírica de Modos (EMD) al análisis del índice MSCI World y su posterior transformación en representaciones gráficas para facilitar el modelado mediante redes neuronales gráficas (GNN). La aplicación de CEEMDAN al índice MSCI World genera nueve funciones de modo intrínseco (IMFs) que capturan componentes oscilatorios en múltiples escalas temporales, desde fluctuaciones de alta frecuencia hasta tendencias de largo plazo. Posteriormente, se implementan cuatro algoritmos de transformación de series a grafos —visibilidad natural, visibilidad horizontal, gráficos de recurrencia y gráficos de transición— para convertir cada IMF en un único grafo. El análisis de las propiedades topológicas revela patrones distintivos según la escala temporal: las IMFs de alta frecuencia generan grafos densos con alta conectividad y estructura de mundo pequeño, mientras que las IMFs de baja frecuencia producen redes más dispersas con mayor longitud de camino característica. Los algoritmos de visibilidad demuestran mayor sensibilidad a la variabilidad de amplitud, generando redes con coeficientes de agrupamiento superiores, mientras que los métodos basados en recurrencia capturan mejor la estructura temporal subyacente. Estos hallazgos establecen las bases para el diseño de arquitecturas GNN adaptadas a las características estructurales específicas de cada componente descompuesto, optimizando el modelado predictivo de series temporales financieras.
\end{abstract}

\vspace{0.5cm}

{\normalfont\normalsize\noindent\textbf{Keywords:} Descomposición Empírica de Modos, EMD, CEEMDAN, Funciones de Modo Intrínseco, MSCI World, Transformación Series-Grafo, Grafos de Visibilidad, Grafos de Recurrencia, Redes Neuronales Gráficas, Series Temporales Financieras, Análisis Topológico}

\newpage
\tableofcontents
\newpage

\section{Introducción}

Analizar series temporales del sector financiero plantea retos específicos derivados del mercado, cuya dinámica no suele ser lineal ni estacionaria. Los métodos tradicionales de descomposición de series temporales, incluidos la transformada wavelet y el análisis de Fourier, se basan en funciones básicas definidas previamente que a menudo resultan inadecuadas para captar las propiedades propias de los conjuntos de datos financieros. La \textbf{Descomposición Empírica de Modos} (EMD), propuesta originalmente por \citep{huang1998empirical}, ha demostrado ser muy adecuada para el análisis de series temporales financieras, ya que proporciona una metodología adaptativa y basada en datos que funciona sin hacer suposiciones sobre la estructura subyacente de los datos. A diferencia de las técnicas tradicionales, la EMD localiza de forma adaptativa los extremos locales y construye envolventes para aislar los componentes oscilatorios, lo que la hace muy eficaz para descomponer series financieras. Cada una de las series resultantes se caracteriza por valores diferentes de amplitud y frecuencia dependientes del tiempo.

La eficacia de la EMD en el análisis de series temporales financieras proviene de la capacidad de esta técnica para adaptarse a las propiedades fundamentales que definen a los datos bursátiles: la \textbf{no estacionariedad} resultante de las tendencias prolongadas, la \textbf{dinámica no lineal} evidente a través de la agrupación de la volatilidad y las transiciones de régimen, y los \textbf{patrones de amplitud y frecuencia variables} según los diferentes estados del mercado. Estas características hacen que las series temporales financieras sean especialmente adecuadas para el análisis EMD, dado que la técnica está diseñada explícitamente para aislar las funciones de modo intrínseco (IMF) de las señales que muestran estas características. El índice \textbf{MSCI World}, que sirve como indicador global del mercado, cumple con las propiedades enumeradas anteriormente, lo que lo posiciona como un caso de estudio adecuado para aplicar EMD en la descomposición de series temporales de este sector.

Aunque el EMD ofrece un marco de descomposición robusto, las IMF resultantes requieren de técnicas de modelización avanzadas para representar relaciones estructurales y temporales. Para cumplir con este objetivo, en esta línea de trabajo se proponen las \textbf{redes neuronales gráficas} (GNN), dada su capacidad para aprender eficazmente a partir de datos estructurados en grafos e identificar patrones tanto localizados como globales a través de mecanismos de paso de mensajes. No obstante, para utilizar las GNN en el análisis de las IMF, es esencial representar las secuencias temporales como nodos y enlaces mediante transformaciones en los datos, manteniendo en todo momento el orden temporal de la serie.

Para sentar las bases del modelado basado en GNN, este artículo se centra en la aplicación de técnicas EMD y la transformación de las componentes en grafos aplicando varios métodos de representación. A un nivel mayor de detalle, se define el alcance de esta investigación que incluye: \textbf{validar las condiciones de aplicabilidad del EMD} para el índice MSCI World mediante pruebas estadísticas; \textbf{descomponer el índice} en más de un elemento utilizando técnicas EMD, con una caracterización exhaustiva de los atributos de energía, varianza, frecuencia y amplitud de cada componente; \textbf{convertir todas las IMF} en representaciones gráficas utilizando algoritmos ya existentes; y \textbf{calcular métricas de los grafos obtenidos} (coeficientes de agrupamiento, medidas de centralidad, patrones de conectividad) que ilustran las características estructurales de cada componente descompuesto. En la ejecución de todos estos pasos se acompaña de una interpretación y discusión de los resultados conseguidos. 

El resto del artículo se organiza en cuatro secciones principales. En la sección~\ref{sec:methodology} se desarrollan los fundamentos teóricos y metodológicos del estudio, donde se describe el conjunto de datos del índice MSCI World, se expone la teoría de la Descomposición Empírica de Modos (EMD) junto con sus variantes, se discuten las consideraciones previas necesarias para su aplicación efectiva, y se presentan los métodos de transformación de series temporales a grafos mediante algoritmos de visibilidad natural (NVG), visibilidad horizontal (HVG) y recurrencia. Posteriormente, la sección~\ref{sec:results} recoge los hallazgos empíricos, los cuales se estructuran en tres bloques: por un lado, la validación de las condiciones adecuadas para la aplicación de EMD (ver sección~\ref{sec:suitcondemd}); por otro, la descomposición EMD del índice MSCI World junto con la caracterización de las funciones modales intrínsecas (IMFs) obtenidas; y finalmente, el análisis de las transformaciones a grafos, abordando sus propiedades estructurales y las métricas calculadas. Para concluir, la sección~\ref{conclfutu} ofrece una síntesis de los principales hallazgos y plantea direcciones para futuras investigaciones.

\section{Metodología}
\label{sec:methodology}

\subsection{Descripción de los datos}

El análisis realizado parte originalmente de los precios de cierre diarios del índice MSCI World~\ref{fig:msci}. El conjunto de datos analizado abarca desde el 12 de enero de 2012 hasta el 26 de noviembre de 2025, y comprende 3490 observaciones diarias. Este periodo cubre diversos regímenes de mercado, incluidos mercados alcistas, correcciones y periodos de alta volatilidad. La serie de precios oscila entre un mínimo de 38,04\$ en dólares americanos (5 de junio de 2012) y un máximo de 186,64\$ (28 de octubre de 2025), lo que representa una rentabilidad acumulada del 376,27 \% durante el periodo de análisis. Los rendimientos diarios se calculan como variaciones porcentuales, con un rendimiento diario medio del 0,0506 \% y una desviación estándar del 1,0780 \%. La volatilidad media móvil de 30 días es del 0,9417 \%, con un valor máximo del 5,1228 \% durante los periodos de alta volatilidad.

\begin{figure}[ht!]
    \centering
    \begin{tikzpicture}
        \node[anchor=south west,inner sep=0] (image) at (0,0) {\includegraphics[width=\textwidth]{./images/english/msci_world_data.png}};
        \begin{scope}[x={(image.south east)},y={(image.north west)}]
            % Etiqueta a) en el panel superior - esquina superior izquierda
            \node[anchor=north west, font=\bfseries, fill=white, inner sep=2pt, opacity=0.8] at (-0.01,0.98) {(a)};
            % Etiqueta b) en el panel medio - esquina superior izquierda
            \node[anchor=north west, font=\bfseries, fill=white, inner sep=2pt, opacity=0.8] at (-0.01,0.65) {(b)};
            % Etiqueta c) en el panel inferior - esquina superior izquierda
            \node[anchor=north west, font=\bfseries, fill=white, inner sep=2pt, opacity=0.8] at (-0.01,0.32) {(c)};
        \end{scope}
    \end{tikzpicture}
    \caption{Datos analizados del MSCI World. \textbf{a)} Evolución de la serie temporal de precios de cierre del índice MSCI World desde enero de 2012 hasta noviembre de 2025, acompañada de las medias móviles de 50 y 200 días que permiten identificar tendencias a corto y largo plazo. \textbf{b)} Retornos diarios calculados como variaciones porcentuales, que capturan la variabilidad diaria del índice.  \textbf{c)} Volatilidad móvil a 30 días, que mide el grado de variación de los retornos en una ventana móvil de 30 días. Esta métrica permite identificar periodos de mayor riesgo y estabilidad del mercado, mostrando picos pronunciados durante eventos de mercado como los observados en 2020.}
    \label{fig:msci} % Para referenciarla en el texto con \ref{fig:mi_imagen}
\end{figure}

\subsection{Teoría y variantes de la Descomposición Empírica de Modos}

La descomposición EMD, introducida por \citep{huang1998empirical}, es un método adaptativo basado en datos para descomponer series temporales no estacionarias y no lineales en un conjunto finito de componentes oscilatorios denominados funciones modales intrínsecas (IMF). A diferencia de los métodos de descomposición tradicionales que se basan en funciones básicas predefinidas (por ejemplo, transformadas de Fourier que utilizan bases sinusoidales o transformadas wavelet que utilizan wavelets madre predefinidas), la EMD identifica de forma adaptativa los extremos locales y construye envolventes para extraer componentes oscilatorios en diferentes escalas temporales.

El principio fundamental de la EMD es el proceso de cribado, que extrae iterativamente las IMF identificando los máximos y mínimos locales, construyendo envolventes superiores e inferiores mediante interpolación cúbica spline, calculando la envolvente media y restándola de la señal. Este proceso continúa hasta que la señal resultante satisface las condiciones de la IMF: el número de extremos y cruces por cero debe diferir en uno como máximo, y el valor medio de las envolventes superior e inferior debe ser cero en cualquier punto. El proceso de cribado finaliza cuando se cumple un criterio de parada, normalmente basado en un umbral de desviación estándar entre iteraciones de cribado consecutivas.

Sin embargo, el EMD estándar posee varias limitaciones, entre las que destaca el problema en la mezcla de modos. Cuando descomponemos con el método EMD original, las oscilaciones de diferentes escalas tienden a mezclarse en un único IMF y las oscilaciones similares pueden dividire entre varios IMF. Este problema se agrava ante la presencia de señales intermitentes y de ruido, que pueden hacer que la descomposición sea inestable y no única.

Para abordar estas limitaciones, se han desarrollado varias variantes de EMD. La \textbf{descomposición empírica de modos en conjunto (EEMD)} \citep{wu2009ensemble} mitiga la mezcla de modos añadiendo ruido blanco a la señal varias veces, realizando EMD en cada versión perturbada por el ruido y promediando los IMF resultantes. El promedio del conjunto cancela el ruido y conserva los componentes reales de la señal, lo que da lugar a descomposiciones más estables y físicamente significativas. La EEMD requiere dos parámetros clave: el número de ensayos del conjunto y la amplitud del ruido, que suele expresarse como una fracción de la desviación estándar de la señal.

En trabajos posteriores, la \textbf{Descomposición empírica modal completa del conjunto con ruido adaptativo (CEEMDAN)} \citep{torres2011complete}, mejora a la EEMD al añadir ruido de forma adaptativa en cada etapa de la descomposición, en lugar de solo en la etapa inicial. Este enfoque reduce el número de ensayos del conjunto necesarios, al tiempo que mantiene la calidad de la descomposición y reduce aún más el ruido residual en las IMF.

Para la obtención de buenos resultados en la descomposición del MSCI World basta con seleccionar la variante EEMD (ver sección \ref{sec:results}) con los parámetros de ejecución definidos en la tabla~\ref{tab:eemd_params}.

\begin{table}[h]
\centering
\caption{Parámetro de la descomposición EEMD}
\label{tab:eemd_params}
\begin{tabular}{ll}
\toprule
\textbf{Parameter} & \textbf{Value} \\
\midrule
Maximum number of IMFs & 14 \\
Standard deviation threshold (SD\_thresh) & 0.25 \\
Minimum sifting iterations (S\_number) & 8 \\
Fixed iterations (FIXE\_H) & 5 \\
Ensemble trials & 100 \\
Noise width & 0.05 (5\% of signal's std dev) \\
\bottomrule
\end{tabular}
\end{table}

\subsection{Consideraciones previas a cumplir para un EMD efectivo.}

Para que el análisis via EMD resulte efectivo, la serie temporal debe cumplir loas siguientes condiciones clave:

\begin{itemize}
\item \textbf{No estacionariedad}: La descomposición por EMD ofrece un funcionamiento óptimo con series no estacionarias, ya que puede capturar de manera adaptativa tendencias y componentes de diferentes escalas temporales. Por ello, se comprueba la estacionariedad del intervalo temporal del MSCI World estudiado utilizando la prueba de Dickey-Fuller aumentada (ADF) \citep{dickey1979distribution} y la prueba de Kwiatkowski-Phillips-Schmidt-Shin (KPSS). \citep{kwiatkowski1992testing}.

\item \textbf{No linearidad}: Aplicar técnicas EMD a la serie no requiere supuestos de linealidad. Se evalua la no linealidad mediante múltiples enfoques: mediante la prueba de Brock-Dechert-Scheinkman (BDS) \citep{brock1996test} sobre los residuos de un modelo AR(1), mediante el análisis de autocorrelación de los rendimientos al cuadrado para detectar efectos ARCH/GARCH \citep{engle1982autoregressive,bollerslev1986generalized}, y finalmente mediante la prueba de Ljung-Box \citep{ljung1978measure} sobre el cuadrado de los rendimientos del índice.

\item \textbf{Amplitud y frecuencia variables}: EMD destaca en la captura de modos con amplitud y frecuencia variables en el tiempo. Por tanto, se analiza la variabilidad de la amplitud mediante medidas de volatilidad continua (ventanas de 252 días correspondientes a un año bursátil) y la variabilidad de la frecuencia mediante análisis de cruce por cero.
\end{itemize}

Los resultados de todos los test listados en este apartado se describen en la sección \ref{sec:suitcondemd}.

\subsection{Transformaciones de series a grafos}

La transformación de series temporales en representaciones gráficas nos permite analizar las propiedades estructurales de los datos temporales utilizando técnicas de teoría de grafos y análisis de redes. Además, convertirlo a grafo habilita la aplicación de redes neuronales gráficas y ofrece una perspectiva alternativa para comprender la dinámica subyacente de los componentes descompuestos de las series temporales. En este estudio, empleamos tres métodos complementarios de transformación de grafos para convertir los IMF extraídos en estructuras de grafos: el grafo de visibilidad natural (NVG), el grafo de visibilidad horizontal (HVG) y las redes basadas en recurrencia.

\subsubsection{Algoritmos de visibilidad}

Los algoritmos de grafos de visibilidad, introducidos por \citep{lacasa2008visibility}, transforman una serie temporal en un grafo estableciendo conexiones entre puntos de datos basadas en criterios de visibilidad geométrica. El principio fundamental es que dos puntos de la serie temporal están conectados si ningún punto intermedio obstruye la línea de visión entre ellos.

La variante denominada algoritmo visibilidad natural (NVG) construye un gráfico en el que cada punto de datos $(t_i, y_i)$ de la serie temporal se convierte en un nodo, y existe un borde entre los nodos $i$ y $j$ si, para todos los puntos intermedios $k$ tales que $i < k < j$, se cumple la siguiente condición:
\begin{equation}
y_k < y_i + (y_j - y_i) \frac{t_k - t_i}{t_j - t_i}
\end{equation}

Este criterio geométrico garantiza que el segmento de línea que conecta los puntos $i$ y $j$ se encuentre por encima de todos los puntos de datos intermedios, estableciendo una relación de visibilidad natural. NVG conserva propiedades importantes de la serie temporal original, como la periodicidad, la fractalidad y la estructura jerárquica, lo que lo hace adecuado para analizar patrones temporales complejos en datos financieros.

En una segunda variante el algoritmo de visibilidad horizontal (HVG) \citep{luque2009horizontal} simplifica a NVG ya que impone una restricción de visibilidad horizontal más estricta. Dos nodos $i$ y $j$ están conectados si, para todos los puntos intermedios $k$ con $i < k < j$, se cumple la siguiente condición:
\begin{equation}
y_k < \min(y_i, y_j)
\end{equation}

Estas dos transformaciones han demostrado tener utilidad en diversas aplicaciones recientes en el ambito financiero. Un ejemplo se tiene en Han et al. (2022) donde desarrollaron MOHVGCA (Multiscale Online-Horizontal-Visibility-Graph Correlation Analysis), un método que emplea HVG para cuantificar correlaciones dependientes de escala entre series temporales financieras, revelando asociaciones dinámicas entre futuros de petróleo WTI y diversas monedas durante crisis petroleras \citep{han2022multiscale}.  Más recientemente, Chen (2025) propuso el uso de grafos de visibilidad ascendente (Rising Visibility Graphs) para extraer motivos topológicos de estructuras de tendencias en índices bursátiles, demostrando que estas estructuras topológicas recurrentes pueden identificarse y utilizarse para pronosticar patrones emergentes en series temporales financieras \citep{chen2025identifying}. Estos estudios evidencian la capacidad de los algoritmos de visibilidad para revelar características estructurales ocultas en series temporales de diversos dominios, validando su aplicación en el análisis de datos financieros complejos como los abordados en este trabajo.

\subsubsection{Algoritmos basados en recurrencia}

Como alternativa a los algoritmos de visibilidad, la matriz de recurrencia puede emplearse como matriz de adyacencia para transformar series temporales en representaciones de grafo mediante la identificación de patrones de recurrencia dentro del sistema dinámico. Este método captura la estructura de recurrencia de la serie en un espacio de embedding de mayor dimensión. Con ello, se consigue desvelar eficazmente tanto la geometría del atractor subyacente como las propiedades dinámicas esenciales del sistema.

La construcción de grafos de recurrencia comienza con la determinación de los parámetros óptimos de embedding utilizando criterios geométricos y de teoría de la información. El parámetro de retardo $\tau$ se selecciona analizando la información mutua, específicamente identificando el primer mínimo de la función de información mutua entre la serie original y su versión retardada en el tiempo. Esto asegura que las coordenadas retardadas proporcionen la máxima información sobre el estado del sistema mientras se minimiza la redundancia. Posteriormente, se estima la dimensión óptima de embedding $d$ utilizando el método de falsos vecinos más cercanos (FNN, por sus siglas en inglés). Esta técnica identifica la dimensión mínima requerida para desplegar el atractor detectando falsos vecinos que aparecen cercanos en dimensiones bajas pero están separados en dimensiones superiores. El método FNN incrementa sistemáticamente la dimensión de embedding hasta que la proporción de falsos vecinos cae por debajo de un umbral, típicamente 10\%.

Una vez determinados $\tau$ (retraso de tiempo) y $d$ (dimensión de embedding), se construye el espacio de embedding utilizando el embedding de coordenadas retardadas. Cada punto $X_i$ en este espacio se define como un vector $d$-dimensional:
\begin{equation}
X_i = [Y_i, Y_{i+\tau}, Y_{i+2\tau}, \ldots, Y_{i+(d-1)\tau}]
\end{equation}
donde $X_i$ representa el $i$-ésimo punto en el espacio de embedding y $Y_i$ denota el valor de la serie temporal original en el tiempo $i$.

El grafo de recurrencia resultante preserva las propiedades topológicas y geométricas del atractor subyacente, permitiendo el análisis de invariantes dinámicos tales como la dimensión de correlación, exponentes de Lyapunov y medidas de entropía. Las características de los nodos en el grafo de recurrencia se asignan utilizando el primer componente de los vectores de embedding, manteniendo una conexión con los valores originales de la serie temporal.

A partir del espacio de embedding generado, el grafo de recurrencia se construye calculando distancias por pares entre todos los vectores de embedding. Dos nodos $i$ y $j$ en el grafo se conectan si la distancia entre sus vectores de embedding correspondientes $X_i$ y $X_j$ está por debajo de un umbral $\varepsilon$. Esto se expresa formalmente a través del elemento $(i, j)$ de la matriz de recurrencia $R$:
\begin{equation}
R_{ij} = \begin{cases}
1 & \text{si } \left\|X_i - X_j\right\| \le \varepsilon \\
0 & \text{en caso contrario}
\end{cases}
\end{equation}
donde $R_{ij}$ es el elemento $(i, j)$-ésimo de la matriz de recurrencia. El umbral $\varepsilon$ se elige típicamente como un percentil de la distribución de distancias por pares, lo que asegura que el grafo capture las estructuras recurrentes más significativas mientras mantiene la tratabilidad computacional.

Se ha demostrado en numerosos estudios que la transformación de series temporales a grafos de recurrencia ha demostrado ser particularmente útil en el análisis de mercados financieros, donde la captura de dinámicas no lineales y patrones complejos es crucial para la predicción y detección de anomalías. En este contexto, Su et al. (2025) \cite{su2025hybrid} desarrollaron un modelo híbrido de doble rama que integra plots de recurrencia con transformers transpuestos para la predicción de tendencias en acciones. En su enfoque, las series temporales de precios de cierre se transforman en plots de recurrencia mediante ventanas deslizantes de segmentos históricos. La primera rama del modelo genera imágenes de RP a partir de estas secuencias univariadas y extrae medidas de análisis de cuantificación de recurrencia (RQA), que capturan características no lineales y relaciones temporales entre puntos. Estas medidas RQA incluyen métricas como la tasa de recurrencia (RR), determinismo (DET) y laminarity (LAM), que cuantifican la predictibilidad y estructura del sistema dinámico. La segunda rama procesa características multivariadas de series temporales utilizando un transformer transpuesto, y ambas ramas se fusionan para clasificar la tendencia del día siguiente, ya sea subida o bajada. Los autores evaluaron el modelo en siete acciones seleccionadas y demostraron que la combinación de características basadas en recurrencia con el modelado de secuencias mediante transformers produce una precisión y F1-score superiores en comparación con métodos tradicionales de aprendizaje automático y deep learning. Este enfoque ilustra cómo los plots de recurrencia pueden capturar dinámicas no lineales y relaciones entre puntos temporales que son difíciles de detectar mediante métodos convencionales de análisis de series temporales.

Por otro lado, Song y Li (2024) \cite{song2024early} propusieron el uso de redes de recurrencia multiplex (MRNs) para la detección de señales de advertencia temprana de transiciones críticas y crashes en mercados bursátiles. En su metodología, vectores multidimensionales de retornos se transforman en redes de recurrencia multiplex, donde cada dimensión, correspondiente a diferentes activos o indicadores, forma una capa de la red. Las relaciones de recurrencia se establecen dentro de cada capa y entre capas, capturando tanto la dinámica individual de cada activo como las interdependencias entre ellos. Los autores calculan indicadores topológicos de las MRNs, específicamente la información mutua promedio y el solapamiento promedio de aristas, para monitorear cambios en la estructura de recurrencia a lo largo del tiempo. Estos indicadores se evalúan tanto en sistemas simulados como en datos empíricos de los mercados de China y Estados Unidos. Los resultados demuestran que los indicadores derivados de MRNs pueden señalar transiciones inminentes de estado en el sistema y proporcionar insights sobre eventos históricos significativos del mercado. La información mutua promedio, en particular, muestra un potencial prometedor como señal de advertencia temprana efectiva. La ventaja principal de las MRNs radica en su capacidad para capturar la estructura de recurrencia multiplex (multidimensional) y los cambios topológicos que las medidas de recurrencia univariadas estándar pueden pasar por alto. Esto permite comparaciones entre mercados y una evaluación más robusta de la efectividad de los indicadores de advertencia temprana. Este enfoque demuestra cómo la extensión de grafos de recurrencia a estructuras multiplex puede mejorar significativamente la capacidad de detectar comportamientos emergentes y transiciones críticas en sistemas financieros complejos.

Tras la revisión de los tipos de transformación, se concluye que la aplicación de los algoritmos de visibilidad y recurrencia a los componentes IMF (Modos Intrínsecos de Función) es de considerable interés. Esta estrategia metodológica permite un análisis más completo del grafo resultante, ya que la visibilidad complementa la información relativa a la estructura temporal de la secuencia, mientras que la recurrencia revela las relaciones del sistema dinámico inherente a la serie temporal. Tal y como se comenta en la sección \ref{conclfutu}, se valorará en próximas líneas de investigación aplicar otro tipo de transformaciones.

\section{Resultados}
\label{sec:results}

\subsection{Condiciones adecuadas para la aplicación de EMD.}
\label{sec:suitcondemd}

Se realizan una batería de pruebas para confirmar un estudio adecuado via descomposición EMD de la serie indexada. Nuestra evaluación exhaustiva del índice MSCI World confirma que la serie temporal cumple las tres condiciones fundamentales para un análisis EMD eficaz. A continuación se presentan los resultados de las tres condiciones evaluadas.

\subsubsection{Estacionariedad}
Las pruebas de estacionariedad proporcionan pruebas sólidas de la no estacionariedad de la serie de precios, lo que es ideal para el EMD. La prueba de Augmented Dickey-Fuller (ADF) sobre la serie de precios de cierre arroja una estadística de prueba de 1.0542 con un valor p de 0.9948, lo que no permite rechazar la hipótesis nula de una raíz unitaria en todos los niveles de significación convencionales. Los valores críticos en los niveles del 1\%, 5\% y 10\% son -3.4322, -2.8624 y -2.5672, respectivamente, todos ellos sustancialmente más negativos que la estadística de prueba. La prueba de Kwiatkowski-Phillips-Schmidt-Shin (KPSS), que tiene la estacionariedad como hipótesis nula, produce una estadística de prueba de 1.2110 con un valor p de 0.0100, lo que rechaza la hipótesis nula de estacionariedad al nivel de significación del 5\%. Por el contrario, la serie de rendimientos con primera diferencia se confirma como estacionaria en ambas pruebas: estadística ADF de -19.5227 (valor p $<$ 0.0001) y estadística KPSS de 0,0341 (valor p = 0.1000). Estos resultados complementarios confirman que la serie de precios muestra un comportamiento no estacionario impulsado por tendencias a largo plazo, lo que la hace muy adecuada para la descomposición EMD.

\subsubsection{Linealidad}
Múltiples pruebas proporcionan pruebas consistentes de una fuerte dinámica no lineal en el índice MSCI World. La prueba de Brock-Dechert-Scheinkman (BDS), aplicada a los residuos de un modelo AR(1), detecta dependencias no lineales significativas en múltiples dimensiones de incrustación. En las dimensiones 3, 4 y 5, las estadísticas de la prueba BDS son 20.99, 24.95 y 27.83, respectivamente, todas con valores p $<$ 0.0001, lo que rechaza firmemente la hipótesis nula de independencia e indica la presencia de una estructura no lineal.

El análisis distributivo respalda aún más la caracterización no lineal. La serie de precios muestra desviaciones significativas de la normalidad, con una asimetría de 0.7515 y una curtosis de -0.3538. La estadística de la prueba de Jarque-Bera arroja un valor p $<$ 0.0001, lo que rechaza firmemente la hipótesis nula de normalidad. Estos hallazgos demuestran colectivamente que el índice MSCI World presenta una rica dinámica no lineal, lo que convierte al EMD en un método de descomposición adecuado.

\subsubsection{Volatilidad en amplitud y frecuencia}

El análisis de autocorrelación de los rendimientos al cuadrado revela una agrupación sustancial de la volatilidad, un rasgo distintivo de los efectos ARCH/GARCH no lineales. Observamos 20 retrasos de autocorrelación significativos de los 20 analizados, con una autocorrelación máxima de 0.9978, lo que indica una fuerte persistencia de la volatilidad. La prueba de Ljung-Box sobre los rendimientos al cuadrado produce una estadística de prueba de 34 210.69 con un valor p $<$ 0.0001, lo que proporciona una evidencia abrumadora de los efectos ARCH/GARCH y confirma la dependencia no lineal en el proceso de varianza.

\subsection{Descomposición EMD}

La descomposición del índice MSCI World se realizó utilizando la variante EEMD (Ensemble Empirical Mode Decomposition) con los parámetros especificados en la tabla~\ref{tab:eemd_params}. Esta elección metodológica se fundamenta en la capacidad de EEMD para mitigar el problema de mezcla de modos mediante la adición de ruido blanco gaussiano y el promediado de múltiples descomposiciones, lo que resulta en una descomposición más estable y físicamente significativa.

La aplicación de EEMD sobre la serie de precios de cierre del índice MSCI World, que comprende 3490 observaciones diarias desde el 12 de enero de 2012 hasta el 26 de noviembre de 2025, produjo una descomposición en \textbf{10 funciones modales intrínsecas (IMFs)} y un \textbf{residuo monótono}. La verificación del residuo confirma que es estrictamente monótono (sin extremos locales), cumpliendo así con el criterio fundamental de terminación del proceso de cribado en EMD.

\subsubsection{Caracterización de las IMFs}

Cada IMF extraída representa oscilaciones en diferentes escalas temporales, desde componentes de alta frecuencia (IMF$_1$) hasta componentes de baja frecuencia (IMF$_{10}$). La tabla~\ref{tab:imf_metrics} presenta un resumen exhaustivo de las métricas clave calculadas para cada IMF: energía, varianza, frecuencia dominante (expresada en ciclos a lo largo del período completo), amplitud promedio y desviación estándar.

\begin{table}[h]
\centering
\caption{Métricas principales de las IMFs extraídas mediante EEMD}
\label{tab:imf_metrics}
\small
\begin{tabular}{lcccc}
\toprule
\textbf{IMF} & \textbf{Energía} & \textbf{Varianza} & \textbf{Frecuencia} & \textbf{Amplitud} \\
 & & & \textbf{(ciclos)} & \textbf{Promedio} \\
\midrule
IMF$_1$ & $1.75 \times 10^3$ & 0.5024 & 1421.0 & 0.8659 \\
IMF$_2$ & $8.39 \times 10^2$ & 0.2405 & 552.0 & 0.5878 \\
IMF$_3$ & $8.74 \times 10^2$ & 0.2505 & 284.0 & 0.5718 \\
IMF$_4$ & $1.55 \times 10^3$ & 0.4430 & 98.0 & 0.7310 \\
IMF$_5$ & $3.27 \times 10^3$ & 0.9359 & 66.0 & 1.0165 \\
IMF$_6$ & $5.78 \times 10^3$ & 1.6523 & 27.0 & 1.3779 \\
IMF$_7$ & $1.31 \times 10^4$ & 3.7452 & 11.0 & 2.0538 \\
IMF$_8$ & $1.31 \times 10^5$ & 37.4858 & 4.0 & 6.5590 \\
IMF$_9$ & $3.71 \times 10^6$ & 153.9538 & 1.0 & 91.9395 \\
IMF$_{10}$ & $2.97 \times 10^7$ & 1058.6908 & 1.0 & 91.9395 \\
\bottomrule
\end{tabular}
\end{table}

El análisis de las métricas revela patrones claros en la estructura de las IMFs. Las IMFs de menor orden (IMF$_1$ a IMF$_4$) exhiben frecuencias dominantes elevadas (de 98 a 1421 ciclos) y amplitudes relativamente pequeñas (menores a 0.87), lo que indica que capturan oscilaciones de corto plazo y ruido de alta frecuencia. A medida que aumenta el orden de la IMF, se observa una disminución sistemática en la frecuencia dominante y un aumento progresivo en la varianza y la amplitud promedio. Las IMFs de mayor orden (IMF$_9$ e IMF$_{10}$) presentan frecuencias dominantes de 1 ciclo y amplitudes promedio elevadas (91.94), sugiriendo que capturan la tendencia de largo plazo de la serie. La IMF$_{10}$ muestra la mayor energía ($2.97 \times 10^7$) y varianza (1058.69), lo que indica que contiene la mayor parte de la variabilidad de la señal original, con una contribución del 88.47\% de la energía total.

\subsubsection{Validación de la descomposición}

La validación de la descomposición EEMD se realizó mediante la reconstrucción de la señal original sumando todas las IMFs. La reconstrucción utilizando únicamente las IMFs (sin incluir el residuo) genera un RMSE de 30.49 y un MAE de 28.50, lo que demuestra la importancia del residuo para capturar la tendencia de largo plazo y asegurar la reconstrucción completa de la señal.

El residuo, calculado como la diferencia entre la señal original y la suma de todas las IMFs, presenta una media de -28.50, una desviación estándar de 27.75, y un rango que va desde 47.01 hasta 132.48. Su naturaleza monótona confirma que el proceso de cribado se completó correctamente y que no quedan oscilaciones adicionales por extraer mediante el algoritmo EEMD. La figura~\ref{fig:imf_decomposition} muestra visualmente la descomposición completa del índice MSCI World, donde se puede observar la progresión desde las IMFs de alta frecuencia (IMF$_1$) hasta las de baja frecuencia (IMF$_{10}$) y el residuo, ilustrando cómo cada componente captura oscilaciones en diferentes escalas temporales y cómo la suma de todos ellos reconstruye la señal original.

La observación de que la suma de las IMFs (sin incluir el residuo) presenta valores sistemáticamente superiores a la señal original, como se evidencia en el análisis de los grafos transformados, es un resultado esperado y válido dentro del marco de la descomposición EEMD. Este comportamiento surge de la naturaleza iterativa del proceso de cribado, donde cada IMF se extrae mediante la eliminación progresiva de envolventes medias, lo que puede introducir pequeñas acumulaciones numéricas que hacen que la suma de las IMFs exceda ligeramente la señal original. En EEMD, el promediado de múltiples realizaciones del conjunto puede amplificar estas discrepancias, ya que cada realización individual puede presentar desviaciones que, al promediarse, no se cancelan completamente. El residuo negativo (media de -28.50) actúa como término de corrección que compensa esta sobreestimación, asegurando que la suma completa (todas las IMFs más el residuo) reconstruya exactamente la señal original. Esta propiedad es consistente con la definición matemática de EEMD, donde la identidad de reconstrucción se cumple estrictamente solo cuando se incluyen tanto las IMFs como el residuo, y no constituye un error del algoritmo sino una característica inherente del proceso de descomposición adaptativa.

\begin{figure}[ht!]
    \centering
    \begin{tikzpicture}
        \node[anchor=south west,inner sep=0] (image) at (0,0) {\includegraphics[width=\textwidth]{./images/english/imf_decomposition.png}};
        \begin{scope}[x={(image.south east)},y={(image.north west)}]
            % Etiqueta a) en el panel superior - esquina superior izquierda
            \node[anchor=north west, font=\bfseries, fill=white, inner sep=2pt, opacity=0.8] at (-0.005,0.98) {(a)};
            % Etiqueta b) en el panel inferior - esquina superior izquierda
            \node[anchor=north west, font=\bfseries, fill=white, inner sep=2pt, opacity=0.8] at (-0.005,0.40) {(b)};
        \end{scope}
    \end{tikzpicture}
    \caption{Descomposición EEMD del índice MSCI World en 10 IMFs y residuo. \textbf{a)} Comparación entre la señal original del índice MSCI World y su reconstrucción utilizando únicamente la suma de las 10 IMFs extraídas mediante EEMD, sin incluir el residuo. \textbf{b)} Error absoluto de reconstrucción, calculado como la diferencia entre la señal original y la suma de las IMFs.}
    \label{fig:imf_decomposition}
\end{figure}

\subsection{Transformación a grafos}

Se realizaron las transformaciones a grafo via algoritmos de visibilidad y de recurrencia para todas las IMF resultantes de la descomposición.

Para ilustrar visualmente la estructura resultante de estas transformaciones, la figura~\ref{fig:hvg_imf} muestra el grafo de visibilidad horizontal (HVG) obtenido a partir de la IMF$_1$, que representa el componente de mayor frecuencia de la descomposición. La visualización revela la topología característica de los grafos HVG: una estructura dispersa donde cada nodo (correspondiente a un punto temporal de la serie) se conecta únicamente con aquellos nodos que cumplen el criterio de visibilidad horizontal. Esta estructura refleja la naturaleza de alta frecuencia de la IMF$_1$, donde las oscilaciones rápidas generan patrones de conectividad localizados que preservan la información temporal de la serie original. La representación gráfica permite apreciar cómo el algoritmo de visibilidad captura las relaciones geométricas entre puntos de datos, transformando la secuencia temporal unidimensional en una estructura de red bidimensional que mantiene las propiedades estructurales del componente modal.

\begin{figure}[ht!]
    \centering
    \includegraphics[width=0.8\textwidth]{./images/english/hvg_imf.png}
    \caption{Grafo de visibilidad horizontal (HVG) resultante de la transformación de la IMF$_1$ del índice MSCI World. Cada nodo representa un punto temporal de la serie, y las conexiones reflejan el criterio de visibilidad horizontal.}
    \label{fig:hvg_imf}
\end{figure}

Específicamente para la construcción de los grafos de recurrencia, se determinaron los parámetros óptimos de embedding de forma independiente para cada IMF mediante los métodos descritos en la sección~\ref{sec:methodology}. El parámetro de retardo $\tau$ se seleccionó mediante el análisis de información mutua, identificando el primer mínimo de la función de información mutua entre la serie original y su versión retardada. La dimensión de embedding $d$ se estimó utilizando el método de falsos vecinos más cercanos (FNN), incrementando sistemáticamente la dimensión hasta que la proporción de falsos vecinos cayera por debajo del umbral del 10\%. El umbral de recurrencia $\varepsilon$ se estableció como el percentil 10 de la distribución de distancias euclidianas por pares entre todos los vectores de embedding, asegurando que el grafo capture las estructuras recurrentes más significativas. Los resultados de la selección de parámetros revelan una variabilidad considerable entre las diferentes IMFs: las IMFs de alta frecuencia (IMF$_1$ a IMF$_4$) requieren dimensiones de embedding de $d=4$ y retardos relativamente pequeños ($\tau$ entre 2 y 6), mientras que las IMFs de baja frecuencia (IMF$_8$ a IMF$_{10}$) y el residuo utilizan dimensiones menores ($d=2$ o $d=3$) y retardos más elevados ($\tau$ entre 16 y 41). Esta variabilidad refleja las diferentes escalas temporales y propiedades dinámicas inherentes a cada componente modal. La tabla~\ref{tab:parametros_recurrencia} presenta los parámetros seleccionados para cada IMF y el residuo.

\begin{table}[h]
\centering
\caption{Parámetros de embedding y umbral $\varepsilon$ por IMF}
\label{tab:parametros_recurrencia}
\small
\begin{tabular}{lccc}
\toprule
\textbf{Componente} & \textbf{$\tau$} & \textbf{$d$} & \textbf{$\varepsilon$} \\
\midrule
IMF$_1$ & 2 & 4 & 0.3620 \\
IMF$_2$ & 2 & 4 & 0.2251 \\
IMF$_3$ & 3 & 4 & 0.1829 \\
IMF$_4$ & 6 & 4 & 0.2123 \\
IMF$_5$ & 11 & 4 & 0.2481 \\
IMF$_6$ & 26 & 3 & 0.1721 \\
IMF$_7$ & 41 & 3 & 0.1324 \\
IMF$_8$ & 33 & 2 & 0.0720 \\
IMF$_9$ & 16 & 2 & 0.0636 \\
IMF$_{10}$ & 25 & 2 & 0.1576 \\
\bottomrule
\end{tabular}
\end{table}

\subsubsection{Propiedades estructurales generales. Conectividad y distancia}

El análisis comparativo de las propiedades estructurales fundamentales de los grafos generados mediante los tres métodos de transformación (HVG, recurrencia y NVG) revela diferencias significativas en la topología de las redes resultantes. Estas diferencias reflejan las características inherentes de cada algoritmo de transformación y su interacción con las propiedades dinámicas de las diferentes IMFs.

En términos de tamaño de red, los tres métodos generan grafos con un número similar de nodos, aproximadamente 3490 nodos correspondientes a las observaciones diarias de la serie temporal. Sin embargo, los grafos de recurrencia presentan ligeramente menos nodos (entre 3408 y 3484) debido al proceso de embedding que requiere parámetros de retardo y dimensión, lo que reduce el número de puntos disponibles para la construcción del grafo. Esta diferencia es mínima y no afecta significativamente el análisis comparativo.

La característica más distintiva entre los métodos de transformación es la densidad de conexiones, medida a través del número de enlaces y la densidad del grafo. Los grafos HVG exhiben una estructura extremadamente dispersa, con densidades que oscilan entre $0.00057$ y $0.00114$, lo que corresponde a entre 3489 y 6968 enlaces. Esta baja densidad refleja la naturaleza restrictiva del criterio de visibilidad horizontal, que solo conecta puntos cuando todos los puntos intermedios están por debajo del mínimo de los dos puntos considerados. En contraste, los grafos de recurrencia presentan densidades moderadas y relativamente estables, con valores entre $0.00264$ y $0.00293$, correspondientes a entre 16033 y 17753 enlaces. Esta densidad intermedia surge del umbral de distancia utilizado en la construcción de la matriz de recurrencia, que captura patrones de similitud en el espacio de embedding.

Los grafos NVG muestran la mayor variabilidad en densidad, con valores que van desde $0.00161$ (para IMF$_1$) hasta $0.70455$ (para IMF$_{10}$), representando un rango de entre 9828 y 4289492 enlaces. Esta variabilidad extrema refleja cómo el criterio de visibilidad natural, menos restrictivo que el HVG, responde de manera diferente a las características de cada IMF. Las IMFs de mayor orden, que capturan tendencias de largo plazo y componentes de baja frecuencia, generan grafos NVG con densidades extremadamente altas, mientras que las IMFs de alta frecuencia producen estructuras más dispersas. Esta propiedad hace que los grafos NVG sean especialmente sensibles a las características espectrales de cada componente modal.

La conectividad global de los grafos, medida a través del número de componentes conexos, revela otra diferencia fundamental entre los métodos. Tanto los grafos HVG como los NVG son siempre completamente conexos (un único componente), lo que facilita el paso de mensajes en arquitecturas GNN y permite el análisis de propiedades globales sin necesidad de técnicas especiales de manejo de componentes desconectados. Por el contrario, los grafos de recurrencia presentan múltiples componentes desconectados, con un número que varía entre 1 (para IMF$_9$, residuo e IMF$_{10}$) y 1024 (para IMF$_3$). Esta fragmentación surge de la naturaleza del espacio de embedding y el umbral de distancia utilizado, donde regiones del espacio de fases que no comparten vecindades cercanas dan lugar a componentes separados. Esta propiedad requiere estrategias especiales en el modelado con GNN, como el procesamiento por componentes o técnicas de pooling que agregan información de múltiples componentes.

El grado promedio de los nodos proporciona una medida adicional de la estructura de conectividad local. Los grafos HVG presentan grados promedio muy bajos, entre $1.99$ y $3.99$, lo que confirma su naturaleza extremadamente dispersa. Los grafos de recurrencia muestran grados promedio moderados y relativamente estables, entre $9.20$ y $9.78$, lo que indica una conectividad local consistente independientemente del orden de la IMF. Los grafos NVG exhiben la mayor variabilidad en grado promedio, desde $5.63$ (para IMF$_1$) hasta $2458.16$ (para IMF$_{10}$), reflejando la dependencia de la densidad del grafo con las características espectrales de cada IMF. Esta variabilidad extrema en los grafos NVG puede plantear desafíos para el aprendizaje de representaciones consistentes en arquitecturas GNN, ya que diferentes IMFs generan estructuras con propiedades de conectividad radicalmente diferentes.

Resumiendo, las tres transformaciones generan estructuras de grafos con propiedades topológicas distintivas. Por un lado los grafos HVG son extremadamente dispersos y siempre conexos mientras que los grafos de recurrencia presentan densidades moderadas pero múltiples componentes desconectados. Por último, los grafos NVG muestran densidades altamente variables que dependen fuertemente del orden de la IMF. Estas diferencias estructurales tienen implicaciones importantes para el diseño de arquitecturas GNN y la selección de estrategias de modelado apropiadas para cada tipo de transformación. La tabla~\ref{tab:propiedades_estructurales} presenta un resumen de todas las propiedades estructurales obtenidas sobre los grafos tras su transformación.

\begin{table}[h]
\centering
\caption{Resumen de propiedades estructurales de los grafos transformados}
\label{tab:propiedades_estructurales}
\small
\begin{tabular}{lccc}
\toprule
\textbf{Propiedad} & \textbf{HVG} & \textbf{Recurrencia} & \textbf{NVG} \\
\midrule
Densidad (rango) & $0.00057$ - $0.00114$ & $0.00264$ - $0.00293$ & $0.00161$ - $0.70455$ \\
Número de enlaces (rango) & $3489$ - $6968$ & $16033$ - $17753$ & $9828$ - $4289492$ \\
Componentes conexos & $1$ (siempre) & $1$ - $1024$ & $1$ (siempre) \\
Diámetro (rango) & $21$ - $3489$ & $20$ - $1110$ & $10$ - $272$ \\
Grado promedio (rango) & $1.99$ - $3.99$ & $9.20$ - $9.78$ & $5.63$ - $2458.16$ \\
Excentricidad promedio (rango) & $16.55$ - $2617.0$ & $13.99$ - $862.25$ & $8.06$ - $266.70$ \\
\bottomrule
\end{tabular}
\end{table}

\subsubsection{Coeficientes de agrupamiento y estructura local}

El coeficiente de agrupamiento, que cuantifica la tendencia de los nodos a formar triángulos o cliques locales, responde a la pregunta de cuan cohesivas resultan las vecindades locales revelando contrastes notables entre los métodos de transformación.

Una comparación directa de los rangos promedio sitúa a los grafos NVG en la cima ($0.5303$-$0.8310$), seguidos por los grafos de recurrencia ($0.2636$-$0.6321$) y finalmente los grafos HVG ($0.0$-$0.6594$). Esta jerarquía refleja diferencias fundamentales en cómo cada criterio de transformación estructura las conexiones locales: el criterio de visibilidad natural (NVG) genera vecindades más densas, mientras que las restricciones geométricas del algoritmo HVG producen estructuras más dispersas.

Sin embargo, esta jerarquía general oculta patrones más sutiles que dependen del orden de la IMF. En los grafos NVG, las IMFs de alta frecuencia (IMF$_1$ a IMF$_4$) mantienen coeficientes entre $0.6678$ y $0.7547$, pero estos valores descienden gradualmente en el rango intermedio (IMF$_5$ a IMF$_9$) hasta alcanzar $0.5303$ para IMF$_9$. La IMF$_{10}$, que captura la tendencia de largo plazo, rompe esta tendencia con el valor más alto ($0.8310$), sugiriendo que los componentes de baja frecuencia generan estructuras locales altamente cohesivas cuando se aplica el criterio de visibilidad natural. Las desviaciones estándar, que oscilan entre $0.1920$ (IMF$_4$) y $0.3332$ (IMF$_9$), revelan una variabilidad moderada dentro de cada grafo. Las medianas, que van de $0.6349$ (IMF$_8$) a $0.8966$ (IMF$_{10}$), confirman que la mayoría de los nodos en los grafos NVG exhiben valores elevados de agrupamiento.

El comportamiento de los grafos HVG contrasta dramáticamente. Aunque las IMFs de alta frecuencia (IMF$_1$ a IMF$_4$) preservan cierta cohesión local ($0.5314$-$0.6594$) a pesar de las restricciones del criterio de visibilidad horizontal, esta cohesión se desvanece completamente en las IMFs de baja frecuencia: IMF$_{10}$ y el residuo alcanzan valores de $0.0$, indicando ausencia total de agrupamiento. Esta degradación sistemática refleja cómo las restricciones geométricas del algoritmo HVG impiden la formación de triángulos en estructuras extremadamente dispersas. La variabilidad, medida por desviaciones estándar entre $0.0$ y $0.2893$, es máxima en las IMFs de alta frecuencia. Las medianas, que varían de $0.0$ a $0.6667$, con valores de $0.5$ para la mayoría de las IMFs intermedias, revelan una distribución bimodal donde muchos nodos carecen completamente de agrupamiento mientras otros mantienen valores moderados.

Los grafos de recurrencia despliegan una dinámica opuesta a los grafos HVG. Mientras que en los HVG la cohesión local disminuye con el orden de la IMF, en los grafos de recurrencia ocurre lo contrario: el residuo ($0.6321$), IMF$_{10}$ ($0.6133$) e IMF$_9$ ($0.4696$) exhiben los valores más altos, mientras que las IMFs de alta frecuencia (IMF$_1$ a IMF$_4$) presentan los más bajos ($0.2636$-$0.3243$). Este patrón inverso sugiere que la estructura de recurrencia captura mejor la cohesión local en componentes de baja frecuencia, donde los patrones dinámicos más estables y recurrentes facilitan la formación de comunidades locales densas. La variabilidad considerable (desviaciones estándar entre $0.0502$ para el residuo y $0.3250$ para IMF$_8$) indica heterogeneidad en la estructura local. Las medianas, que oscilan entre $0.3333$ (IMF$_3$ e IMF$_4$) y $0.6429$ (IMF$_{10}$), reflejan esta variabilidad según el orden de la IMF.

La tabla~\ref{tab:clustering_coefficients} sintetiza estos resultados, mostrando los rangos de valores promedio, medianos y desviaciones estándar para cada método.

\begin{table}[h]
\centering
\caption{Coeficientes de agrupamiento por método de transformación y orden de IMF}
\label{tab:clustering_coefficients}
\small
\begin{tabular}{lccc}
\toprule
\textbf{Método} & \textbf{Coeficiente promedio} & \textbf{Coeficiente mediano} & \textbf{Desviación estándar} \\
 & \textbf{(rango)} & \textbf{(rango)} & \textbf{(rango)} \\
\midrule
HVG & $0.0$ - $0.6594$ & $0.0$ - $0.6667$ & $0.0$ - $0.2893$ \\
Recurrencia & $0.2636$ - $0.6321$ & $0.3333$ - $0.6429$ & $0.0502$ - $0.3250$ \\
NVG & $0.5303$ - $0.8310$ & $0.6349$ - $0.8966$ & $0.1920$ - $0.3332$ \\
\bottomrule
\end{tabular}
\end{table}

La estructura local de los grafos transformados, caracterizada a través de los coeficientes de agrupamiento, revela propiedades topológicas distintivas que tienen implicaciones importantes para el modelado con GNN. Los grafos NVG, con sus altos coeficientes de agrupamiento, presentan vecindades locales densas y cohesivas que facilitan el paso de mensajes localizado y pueden capturar patrones de interacción complejos. Los grafos HVG, con sus coeficientes de agrupamiento variables y frecuentemente bajos, especialmente para IMFs de baja frecuencia, generan estructuras locales más dispersas que pueden requerir estrategias de agregación diferentes en las capas de convolución gráfica. Los grafos de recurrencia, con sus coeficientes de agrupamiento intermedios y su patrón inverso según el orden de la IMF, ofrecen una estructura local que refleja las propiedades dinámicas del sistema subyacente, donde los componentes de baja frecuencia muestran mayor cohesión local debido a la naturaleza recurrente de sus patrones dinámicos.

\subsubsection{Análisis de centralidad}

Identificar qué nodos desempeñan roles críticos en la estructura del grafo resulta fundamental para comprender cómo fluye la información y dónde se concentra la influencia. Para este propósito, examinamos tres métricas de centralidad que capturan aspectos complementarios: la centralidad de intermediación (betweenness) identifica puentes estratégicos, la centralidad de cercanía (closeness) revela nodos con acceso rápido al resto de la red, y la centralidad de vector propio (eigenvector) mide la influencia propagada a través de conexiones.

Cuando se analiza qué nodos actúan como puentes críticos en el flujo de información, emergen contrastes marcados entre los métodos de transformación. Los grafos NVG muestran valores promedio de \textbf{centralidad de intermediación} extremadamente bajos (entre $0.0006$ y $0.0058$), lo que sugiere que la alta densidad distribuye uniformemente el rol de intermediación, evitando la dependencia de nodos específicos como puentes. Por el contrario, los grafos HVG exhiben una variabilidad dramática: mientras que las IMFs de alta frecuencia (IMF$_1$) presentan valores promedio de $0.0026$, las IMFs de baja frecuencia (IMF$_{10}$ y residuo) alcanzan $0.3333$, con valores máximos que llegan hasta $0.6539$ para IMF$_6$. Esta disparidad surge porque la estructura extremadamente dispersa de las IMFs de baja frecuencia concentra el flujo de información en pocos nodos estratégicos. Los grafos de recurrencia ocupan una posición intermedia, con valores promedio entre $0.0006$ (IMF$_3$) y $0.1094$ (IMF$_{10}$), donde la presencia de múltiples componentes desconectados reduce la necesidad de nodos intermediarios centralizados.

La capacidad de acceder rápidamente al resto de la red, medida por la \textbf{centralidad de cercanía}, muestra un patrón distinto. Los grafos NVG destacan con los valores más elevados (rango entre $0.0680$ para IMF$_9$ y $0.2503$ para IMF$_7$), confirmando que su alta densidad permite que la mayoría de los nodos mantengan distancias cortas entre sí. Tanto los grafos HVG como los de recurrencia presentan valores considerablemente menores, aunque con diferencias sutiles: en los grafos HVG, los valores oscilan entre $0.0009$ (IMF$_{10}$ y residuo) y $0.1019$ (IMF$_1$), mientras que en los grafos de recurrencia el rango va de $0.0028$ (IMF$_{10}$) a $0.0993$ (IMF$_1$). En ambos casos, las IMFs de alta frecuencia muestran mayor accesibilidad debido a sus estructuras más compactas y diámetros reducidos.

La \textbf{centralidad de vector propio}, que evalúa la influencia basada en la importancia de los vecinos, revela limitaciones estructurales importantes. Los grafos HVG y de recurrencia comparten valores promedio muy bajos (entre $0.0016$-$0.0058$ y $0.0025$-$0.0058$, respectivamente), aunque con diferencias en los valores máximos: los grafos HVG alcanzan hasta $0.6186$ para IMF$_6$, mientras que los grafos de recurrencia rara vez superan $0.1876$. Esta disparidad en los máximos refleja que, aunque la estructura dispersa de los grafos HVG limita la propagación general de influencia, algunos nodos pueden mantener influencia local significativa. Los grafos NVG presentan valores promedio más elevados (entre $0.0045$ y $0.0159$), con valores máximos que varían considerablemente desde $0.0698$ (IMF$_7$) hasta $0.4336$ (IMF$_2$), sugiriendo que la alta densidad facilita la propagación de influencia, aunque de manera más distribuida en las IMFs de mayor orden.

\begin{table}[h]
\centering
\caption{Medidas de centralidad por método de transformación}
\label{tab:centralidad}
\small
\begin{tabular}{lccc}
\toprule
\textbf{Métrica} & \textbf{HVG} & \textbf{Recurrencia} & \textbf{NVG} \\
 & \textbf{(rango promedio)} & \textbf{(rango promedio)} & \textbf{(rango promedio)} \\
\midrule
Betweenness & $0.0026$ - $0.3333$ & $0.0006$ - $0.1094$ & $0.0006$ - $0.0058$ \\
Closeness & $0.0009$ - $0.1019$ & $0.0028$ - $0.0993$ & $0.0680$ - $0.2503$ \\
Eigenvector & $0.0016$ - $0.0058$ & $0.0025$ - $0.0058$ & $0.0045$ - $0.0159$ \\
\bottomrule
\end{tabular}
\end{table}

La tabla~\ref{tab:centralidad} sintetiza los rangos observados para cada métrica. Estas diferencias estructurales tienen implicaciones directas para el modelado: los grafos NVG facilitan la propagación eficiente de información pero con influencia distribuida, los grafos HVG generan estructuras donde ciertos nodos adquieren roles críticos como intermediarios (especialmente en IMFs de baja frecuencia), mientras que los grafos de recurrencia, con valores generalmente bajos en todas las medidas, reflejan una importancia estructural más localizada dentro de componentes fragmentados.

\subsubsection{Patrones según el orden de las IMFs}

El análisis de las propiedades estructurales de los grafos transformados revela patrones sistemáticos que varían según el orden de las IMFs, reflejando las diferentes escalas temporales y propiedades dinámicas inherentes a cada componente modal. Estos patrones proporcionan insights valiosos sobre cómo las características espectrales de las IMFs se traducen en propiedades topológicas de los grafos resultantes.

Los \textbf{grafos HVG} exhiben una degradación estructural progresiva conforme aumenta el orden de la IMF. En las IMFs de alta frecuencia (IMF$_1$ a IMF$_4$), las estructuras resultan relativamente compactas: diámetros entre $21$ (IMF$_1$) y $63$ (IMF$_4$), densidades moderadas (entre $0.0011$ y $0.0011$), y coeficientes de agrupamiento elevados (entre $0.5314$ y $0.6594$). Sin embargo, esta cohesión inicial se desvanece sistemáticamente en las IMFs de mayor orden. Los diámetros se disparan hasta $3489$ para IMF$_{10}$ y residuo, mientras que las densidades caen a $0.0006$ y los coeficientes de agrupamiento alcanzan cero. La explicación subyacente radica en que el criterio restrictivo de visibilidad horizontal, aunque efectivo para capturar oscilaciones rápidas, resulta cada vez más limitante cuando se aplica a componentes de baja frecuencia, produciendo estructuras extremadamente lineales y dispersas sin cohesión local.

En contraste con los grafos HVG, los \textbf{grafos de recurrencia} muestran una dinámica opuesta: las IMFs de baja frecuencia desarrollan propiedades estructurales notablemente más cohesivas. Las IMFs de alta frecuencia (IMF$_1$ a IMF$_4$) se caracterizan por múltiples componentes desconectados (entre $539$ y $1024$ componentes), diámetros relativamente pequeños (entre $20$ y $34$), y coeficientes de agrupamiento bajos (entre $0.2636$ y $0.3243$). La transición hacia IMFs de mayor orden conlleva una consolidación estructural marcada: el número de componentes se reduce hasta $1$ para IMF$_9$, IMF$_{10}$ y residuo, los diámetros crecen hasta $1110$ para IMF$_9$, y los coeficientes de agrupamiento se incrementan significativamente hasta $0.6321$ para el residuo. Esta evolución se debe a que los componentes de baja frecuencia, con patrones dinámicos más estables y recurrentes, generan estructuras de embedding más cohesivas donde los puntos del espacio de fases se agrupan más densamente, favoreciendo la formación de comunidades locales más compactas.

El comportamiento de los \textbf{grafos NVG} difiere sustancialmente de los anteriores, mostrando variaciones no monótonas en las propiedades estructurales según el orden de la IMF. Las IMFs de alta frecuencia (IMF$_1$ a IMF$_4$) se distinguen por densidades relativamente bajas (entre $0.0016$ y $0.0052$), diámetros pequeños (entre $10$ y $14$), y coeficientes de agrupamiento elevados (entre $0.6678$ y $0.7547$). En el rango intermedio (IMF$_5$ a IMF$_9$), la densidad aumenta progresivamente hasta $0.1703$ para IMF$_9$, los diámetros crecen moderadamente hasta $127$ para IMF$_9$, mientras que los coeficientes de agrupamiento descienden gradualmente hasta $0.5303$ para IMF$_9$. La IMF$_{10}$ rompe esta tendencia con propiedades extremas: densidad de $0.7045$, diámetro de $272$, y el coeficiente de agrupamiento más alto ($0.8310$). Esta variabilidad no monótona surge porque el criterio de visibilidad natural responde de manera heterogénea a las características espectrales de cada IMF, generando estructuras altamente densas en componentes de baja frecuencia donde la mayoría de los puntos satisfacen el criterio de visibilidad.

Estos patrones sistemáticos según el orden de las IMFs tienen implicaciones importantes para el diseño de arquitecturas GNN. Las IMFs de alta frecuencia, con sus estructuras más compactas y cohesivas en grafos HVG y NVG, pueden beneficiarse de capas de convolución gráfica estándar que aprovechan la cohesión local. Las IMFs de baja frecuencia requieren estrategias adaptativas: en grafos HVG, donde la estructura es extremadamente dispersa, pueden necesitarse técnicas de atención o pooling adaptativo; en grafos de recurrencia, donde la estructura es más cohesiva, las capas estándar pueden ser efectivas; y en grafos NVG, donde la densidad es extremadamente alta, pueden requerirse técnicas de muestreo o pooling para manejar la complejidad computacional.

\subsubsection{Implicaciones para el modelado con GNN}

Las propiedades estructurales identificadas en los grafos transformados tienen implicaciones importantes para la aplicación de redes neuronales gráficas. Los grafos NVG, con su alta densidad y estructura a muy corta distancia, proporcionan un entorno rico para el paso de mensajes en GNN, permitiendo que la información se propague eficientemente a través de la red. Sin embargo, la alta densidad también puede aumentar la complejidad computacional y potencialmente introducir ruido en las representaciones aprendidas.

Los grafos HVG, con su estructura más dispersa y regular, ofrecen una representación más parsimoniosa que puede ser computacionalmente más eficiente, pero la variabilidad extrema en las distancias (especialmente para IMFs de baja frecuencia) puede plantear desafíos para el aprendizaje de representaciones consistentes. La estructura conexa de todos los grafos HVG es una ventaja, ya que permite el paso de mensajes global sin necesidad de técnicas especiales para manejar componentes desconectados.

Los grafos de recurrencia presentan el desafío adicional de múltiples componentes desconectados para la mayoría de las IMFs, lo que requiere estrategias especiales en el diseño de GNN, como el procesamiento por componentes o técnicas de pooling que agregan información de múltiples componentes. Sin embargo, la estructura de recurrencia captura propiedades dinámicas fundamentales del sistema que pueden ser valiosas para el modelado predictivo.

La complementariedad de los tres métodos sugiere que una estrategia de modelado híbrida, que combine información de múltiples representaciones gráficas, podría aprovechar las fortalezas de cada método mientras mitiga sus limitaciones individuales. Esta observación sienta las bases para futuras investigaciones en el diseño de arquitecturas GNN multi-vista para el análisis de series temporales financieras descompuestas mediante EMD.


\section{Conclusiones y líneas futuras}
\label{conclfutu}

Como resultado del estudio se tiene el desarrollo de un marco metodológico que combina la Descomposición Empírica de Modos con técnicas de transformación a grafos para analizar series temporales financieras, aplicado específicamente al índice MSCI World. Los resultados obtenidos muestran que esta integración de técnicas complementarias permite una comprensión más profunda de la dinámica compleja, no lineal y no estacionaria que caracteriza a los mercados financieros.

La validación estadística mediante las pruebas ADF, KPSS y BDS confirmó que el índice MSCI World cumple las condiciones necesarias para aplicar EMD de manera efectiva, estableciendo una base sólida para el análisis posterior. La descomposición mediante EEMD produjo 10 Funciones de Modo Intrínseco más un residuo monótono, cada uno capturando diferentes escalas temporales de la dinámica del mercado. Resulta especialmente relevante el hallazgo de que las IMFs de orden superior concentran la mayor parte de la energía y varianza del sistema, donde la IMF$_{10}$ contiene el 88.47\% de la energía total, representando las tendencias de largo plazo, mientras que las IMFs de orden inferior capturan oscilaciones de alta frecuencia y ruido de corto plazo.

La transformación de estas IMFs en representaciones de grafos mediante tres métodos complementarios reveló propiedades topológicas distintivas que reflejan diferentes aspectos de la estructura subyacente de los datos financieros. Los grafos HVG generan estructuras extremadamente dispersas con baja densidad y grado promedio, pero mantienen conectividad completa. Sin embargo, exhiben alta variabilidad en las distancias para IMFs de baja frecuencia, lo que complica el aprendizaje de representaciones consistentes. Las redes de recurrencia muestran densidades moderadas y estables, pero frecuentemente presentan múltiples componentes desconectadas, especialmente para IMFs de alta frecuencia. Esta fragmentación requiere estrategias especializadas de diseño de arquitecturas de redes neuronales para el procesamiento efectivo de componentes. Los grafos NVG exhiben la mayor variabilidad en densidad, desde estructuras dispersas para IMFs de alta frecuencia hasta grafos extremadamente densos para IMFs de baja frecuencia. Esta alta densidad, aunque captura relaciones ricas, incrementa significativamente la complejidad computacional.

Una hallazgo clave obtenido tras el análisis es que las propiedades estructurales de los grafos varían sistemáticamente con el orden de las IMFs, pero cada método de transformación exhibe patrones distintivos. Los HVG muestran degradación estructural progresiva, las redes de recurrencia presentan mayor cohesión para IMFs de baja frecuencia, y los NVG exhiben variaciones no monótonas. Esta diversidad sugiere que ningún método de representación de grafos es óptimo para todos los componentes IMF.

Los resultados de este estudio proporcionan información crucial para el diseño de arquitecturas de GNN aplicadas al pronóstico de series temporales financieras. Las propiedades topológicas identificadas (dispersión extrema en HVG, componentes desconectadas en redes de recurrencia, y alta densidad en NVG) plantean desafíos específicos de modelado que deben ser abordados mediante estrategias arquitectónicas especializadas. En particular, se requieren enfoques híbridos que puedan aprovechar las fortalezas complementarias de múltiples representaciones de grafos simultáneamente.

En conclusión, este trabajo establece que la integración de EMD con análisis de grafos constituye una herramienta poderosa para desentrañar la estructura multinivel de las series temporales financieras. Los patrones topológicos descubiertos no solo profundizan nuestra comprensión de la dinámica del mercado, sino que también informan el desarrollo de sistemas de pronóstico más sofisticados basados en GNN. La metodología propuesta es generalizable a otros índices financieros y clases de activos, abriendo vías prometedoras para la investigación futura en el análisis cuantitativo de mercados financieros.

Los resultados y limitaciones identificados en este trabajo abren múltiples direcciones prometedoras para la investigación futura. La extensión natural de este trabajo es el desarrollo e implementación de arquitecturas GNN especializadas que aprovechen las propiedades topológicas identificadas en las representaciones de grafos de las IMFs. Dado que cada método de transformación a grafos captura aspectos complementarios de la dinámica del mercado, se propone desarrollar arquitecturas GNN multi-vista que procesen simultáneamente múltiples representaciones de grafos \citep{liu2023multimodal, wang2023hatr}. Estas arquitecturas podrían emplear mecanismos de atención adaptativa para ponderar dinámicamente la contribución de cada vista según el horizonte temporal de predicción y las condiciones del mercado \citep{su2024attention}.

La literatura reciente demuestra la efectividad de combinar codificadores temporales con capas de convolución de grafos para capturar tanto dependencias temporales como relaciones espaciales entre activos \citep{liu2023multimodal, wang2023hatr, jin2023temporal}. Para las IMFs extraídas en este estudio, se recomienda desarrollar modelos que adapten dinámicamente la estructura del grafo a través del tiempo, permitiendo que las relaciones entre componentes evolucionen según los regímenes del mercado. Las redes de recurrencia identificadas en este trabajo sugieren la presencia de relaciones de orden superior entre las IMFs. Arquitecturas basadas en hipergrafos adaptativos \citep{su2024attention, agarwal2022think} podrían capturar estas asociaciones multi-componente de manera más efectiva que los grafos tradicionales. Específicamente, se propone investigar redes de atención de hipergrafos en espacios hiperbólicos para modelar relaciones jerárquicas entre IMFs de diferentes escalas temporales \citep{agarwal2022think}, grafos heterogéneos que incorporen múltiples tipos de relaciones entre componentes \citep{jin2023temporal}, y mecanismos de atención dual que operen tanto a nivel de nodo como de tipo de relación para refinar la propagación de información \citep{wang2023hatr}.

La escalabilidad es fundamental para la aplicación práctica de GNNs en mercados financieros, donde el número de activos, la frecuencia de actualización de datos y la complejidad de las relaciones pueden crecer exponencialmente. El manejo de grafos en entornos distribuidos representa una línea de investigación crítica. Se propone investigar la adaptación de frameworks modernos de procesamiento distribuido de grafos para el contexto específico de series temporales financieras. Los sistemas distribuidos contemporáneos emplean estrategias de particionamiento de grafos, paralelismo multi-nivel y optimizaciones de comunicación para escalar el entrenamiento de GNNs \citep{xu2025capturing}. Para este trabajo, se recomienda evaluar estrategias de particionamiento que respeten la estructura temporal de las IMFs, minimizando la comunicación entre trabajadores, implementar paralelismo de datos para procesar múltiples ventanas temporales simultáneamente, y desarrollar esquemas de caché y pre-cómputo de embeddings para IMFs estables de largo plazo.

Los mercados financieros generan datos continuamente, requiriendo actualización incremental de las representaciones de grafos. Se propone investigar algoritmos de actualización incremental de EMD que permitan incorporar nuevas observaciones sin recomputar toda la descomposición, arquitecturas GNN para grafos dinámicos que soporten actualizaciones continuas de estructura y atributos \citep{gravina2023scalable, liu2023decoupled}, y sistemas de inferencia en tiempo real que balanceen latencia y precisión mediante cómputo distribuido. Los grafos NVG densos y las redes de recurrencia con múltiples componentes presentan desafíos de huella de memoria y ancho de banda de comunicación en entornos distribuidos. Líneas de investigación incluyen técnicas de compresión de grafos que preserven propiedades topológicas esenciales para el aprendizaje, estrategias de muestreo de vecindarios adaptativo que reduzcan la complejidad sin degradar la calidad de predicción, y protocolos de comunicación asíncrona para minimizar sincronizaciones globales durante el entrenamiento distribuido.

Extender el análisis de un único índice a portafolios completos de activos globales requiere capacidades de procesamiento distribuido robustas. Se propone desarrollar arquitecturas federadas que permitan entrenar modelos sobre datos de múltiples mercados preservando privacidad, implementar sistemas de transferencia de aprendizaje que aprovechen conocimiento de un mercado para mejorar predicciones en otros, y diseñar estrategias de agregación multi-índice que capturen dependencias entre mercados regionales y globales.

Otras direcciones de investigación incluyen aplicar el marco metodológico propuesto a diferentes instrumentos financieros para evaluar la generalización de los patrones topológicos identificados, investigar cómo las propiedades topológicas de los grafos derivados de IMFs varían durante diferentes regímenes de mercado y explotar estos patrones para detección temprana de cambios de régimen, desarrollar técnicas de interpretación para modelos GNN que permitan identificar qué componentes IMF y qué estructuras de grafo contribuyen más a las predicciones, investigar métodos adaptativos para seleccionar automáticamente parámetros óptimos de EEMD según las características específicas de cada serie temporal, valorar otro tipo de transformaciones de series temporales a grafos y explorar la incorporación de fuentes de datos no tradicionales en las representaciones de grafos para enriquecer el poder predictivo de los modelos GNN.

En resumen, las líneas futuras propuestas abarcan desde el desarrollo de arquitecturas GNN especializadas que aprovechen las propiedades topológicas identificadas, hasta la implementación de sistemas distribuidos escalables capaces de procesar grafos dinámicos en tiempo real. La convergencia de técnicas de descomposición de señales, teoría de grafos, aprendizaje profundo y computación distribuida promete avances significativos en el pronóstico de series temporales financieras, con potencial impacto tanto en la investigación académica como en aplicaciones prácticas de trading y gestión de riesgos.

\begin{thebibliography}{99}

\bibitem{bollerslev1986generalized}
Bollerslev, T. (1986).
\newblock Generalized autoregressive conditional heteroskedasticity.
\newblock {\em Journal of Econometrics}, 31(3), 307--327.

\bibitem{brock1996test}
Brock, W.~A., Dechert, W.~D., Scheinkman, J.~A., \& LeBaron, B. (1996).
\newblock A test for independence based on the correlation dimension.
\newblock {\em Econometric Reviews}, 15(3), 197--235.

\bibitem{chen2025identifying}
Zeng, Z., \& Chen, Y. (2025).
\newblock Identifying and forecasting recurrently emerging stock trend structures via rising visibility graphs.
\newblock {\em Forecasting}, 7(2), 26.
\newblock \url{https://doi.org/10.3390/forecast7020026}

\bibitem{dickey1979distribution}
Dickey, D.~A., \& Fuller, W.~A. (1979).
\newblock Distribution of the estimators for autoregressive time series with a unit root.
\newblock {\em Journal of the American Statistical Association}, 74(366a), 427--431.

\bibitem{eckmann1987recurrence}
Eckmann, J.~P., Kamphorst, S.~O., \& Ruelle, D. (1987).
\newblock Recurrence plots of dynamical systems.
\newblock {\em EPL (Europhysics Letters)}, 4(9), 973.

\bibitem{engle1982autoregressive}
Engle, R.~F. (1982).
\newblock Autoregressive conditional heteroscedasticity with estimates of the variance of United Kingdom inflation.
\newblock {\em Econometrica}, 50(4), 987--1007.

\bibitem{han2022multiscale}
Han, M., Fan, Q., \& Ling, G. (2022).
\newblock Multiscale online-horizontal-visibility-graph correlation analysis of financial market.
\newblock {\em Physica A: Statistical Mechanics and Its Applications}, 607, 128195.
\newblock \url{https://doi.org/10.1016/j.physa.2022.128195}

\bibitem{huang1998empirical}
Huang, N.~E., Shen, Z., Long, S.~R., Wu, M.~C., Shih, H.~H., Zheng, Q., \ldots \& Liu, H.~H. (1998).
\newblock The empirical mode decomposition and the Hilbert spectrum for nonlinear and non-stationary time series analysis.
\newblock {\em Proceedings of the Royal Society of London. Series A: Mathematical, Physical and Engineering Sciences}, 454(1971), 903--995.

\bibitem{kwiatkowski1992testing}
Kwiatkowski, D., Phillips, P.~C., Schmidt, P., \& Shin, Y. (1992).
\newblock Testing the null hypothesis of stationarity against the alternative of a unit root: How sure are we that economic time series have a unit root?
\newblock {\em Journal of Econometrics}, 54(1-3), 159--178.

\bibitem{lacasa2008visibility}
Lacasa, L., Luque, B., Ballesteros, F., Luque, J., \& Nuño, J.~C. (2008).
\newblock From time series to complex networks: The visibility graph.
\newblock {\em Proceedings of the National Academy of Sciences}, 105(13), 4972--4975.

\bibitem{ljung1978measure}
Ljung, G.~M., \& Box, G.~E. (1978).
\newblock On a measure of lack of fit in time series models.
\newblock {\em Biometrika}, 65(2), 297--303.

\bibitem{luque2009horizontal}
Luque, B., Lacasa, L., Ballesteros, F., \& Luque, J. (2009).
\newblock Horizontal visibility graphs: Exact results for random time series.
\newblock {\em Physical Review E}, 80(4), 046103.

\bibitem{song2024early}
Song, S., \& Li, H. (2024).
\newblock Early warning signals for stock market crashes: empirical and analytical insights utilizing nonlinear methods.
\newblock {\em EPJ Data Science}, 13(1), 16.
\newblock \url{https://doi.org/10.1140/epjds/s13688-024-00457-2}

\bibitem{su2025hybrid}
Su, J., Li, H., Wang, R., Guo, W., Hao, Y., Kurths, J., \& Gao, Z. (2025).
\newblock A hybrid dual-branch model with recurrence plots and transposed transformer for stock trend prediction.
\newblock {\em Chaos: An Interdisciplinary Journal of Nonlinear Science}, 35(1), 013125.
\newblock \url{https://doi.org/10.1063/5.0233275}

\bibitem{torres2011complete}
Torres, M.~E., Colominas, M.~A., Schlotthauer, G., \& Flandrin, P. (2011).
\newblock A complete ensemble empirical mode decomposition with adaptive noise.
\newblock {\em 2011 IEEE International Conference on Acoustics, Speech and Signal Processing (ICASSP)}, 4144--4147.

\bibitem{wu2009ensemble}
Wu, Z., \& Huang, N.~E. (2009).\textbf{}
\newblock Ensemble empirical mode decomposition: a noise-assisted data analysis method.
\newblock {\em Advances in Adaptive Data Analysis}, 1(01), 1--41.

\bibitem{liu2023multimodal}
Liu, R., Liu, H., Huang, H., Yang, J., Gao, Y., Zhu, J., Wang, C., \& Xu, K. (2023).
\newblock Multimodal multiscale dynamic graph convolution networks for stock price prediction.
\newblock {\em Pattern Recognition}, 144, 110211.
\newblock \url{https://doi.org/10.1016/j.patcog.2023.110211}

\bibitem{wang2023hatr}
Wang, H., Wang, T., Li, S., Li, J., \& Zhang, Z. (2023).
\newblock HATR-I: Hierarchical Adaptive Temporal Relational Interaction for Stock Trend Prediction.
\newblock {\em IEEE Transactions on Knowledge and Data Engineering}, 35(7), 7492--7506.
\newblock \url{https://doi.org/10.1109/TKDE.2022.3188320}

\bibitem{su2024attention}
Su, H., Wang, X., Qin, Y., Zhao, X., \& Lin, C. (2024).
\newblock Attention based adaptive spatial--temporal hypergraph convolutional networks for stock price trend prediction.
\newblock {\em Expert Systems with Applications}, 237, 121899.
\newblock \url{https://doi.org/10.1016/j.eswa.2023.121899}

\bibitem{jin2023temporal}
Jin, M., Koh, H.~Y., Wen, Q., Zambon, D., Alippi, C., Webb, G.~I., King, I., \& Pan, S. (2023).
\newblock Temporal and Heterogeneous Graph Neural Network for Financial Time Series Prediction.
\newblock In {\em Proceedings of the 31st ACM International Conference on Information and Knowledge Management} (CIKM '22), 3584--3593.
\newblock \url{https://doi.org/10.1145/3511808.3557089}

\bibitem{agarwal2022think}
Agarwal, S., Sawhney, R., Thakkar, M., \& Gupta, D. (2022).
\newblock THINK: Temporal Hypergraph Hyperbolic Network.
\newblock In {\em Proceedings of the 2022 IEEE International Conference on Data Mining} (ICDM), 1--10.
\newblock \url{https://doi.org/10.1109/ICDM54844.2022.00096}

\bibitem{gravina2023scalable}
Gravina, A., Zambon, D., Bacciu, D., \& Alippi, C. (2023).
\newblock Scalable spatiotemporal graph neural networks.
\newblock In {\em Proceedings of the AAAI Conference on Artificial Intelligence}, 37(6), 7688--7696.

\bibitem{liu2023decoupled}
Liu, Y., Li, K., Xie, Y., Guo, Y., Yan, J., \& Yu, P.~S. (2023).
\newblock Decoupled Graph Neural Networks for Large Dynamic Graphs.
\newblock {\em arXiv preprint arXiv:2305.08273}.
\newblock \url{https://doi.org/10.48550/arXiv.2305.08273}

\bibitem{xu2025capturing}
Xu, Q. (2025).
\newblock Capturing Structural Evolution in Financial Markets with Graph Neural Time Series Models.
\newblock {\em Preprints}, 2025071260.
\newblock \url{https://doi.org/10.20944/preprints202507.1260.v1}

\end{thebibliography}

\end{document}
